\usepackage[a4paper,left=2.75cm,right=2.6cm,top=2.5cm,bottom=2.5cm,heightrounded]{geometry}
\usepackage{emptypage}
\usepackage[spanish,es-tabla]{babel}
\unaccentedoperators

\usepackage{csquotes}
\usepackage{graphicx}
\usepackage{wrapfig}
\usepackage{mathtools}
\usepackage{xspace}
\usepackage{xeCJK}
\usepackage{microtype}
% --- 3. MATEMÁTICAS Y TIPOGRAFÍA ---
\usepackage{amsmath,amsthm,amsfonts}
\usepackage{newpxtext,newpxmath} % Tipografía Palatino
\linespread{1.03}

% --- 4. GRÁFICOS Y COLORES ---
\usepackage[dvipsnames]{xcolor}
\definecolor{goban}{HTML}{efce8f}
\definecolor{othello}{HTML}{485d3f}
\usepackage{float}
\usepackage{pdfpages}

\usepackage{tikz}
\usepackage{tikz-cd}
\usepackage{pgfplots}
\pgfplotsset{compat=1.18}
% Todas las librerías de TikZ agrupadas en una sola línea
\usetikzlibrary{babel, arrows.meta, bending, positioning, shadows, calc, shapes.geometric, decorations.text}

% --- 5. TABLAS Y UTILIDADES ---
\usepackage{multicol}
\usepackage{multirow}
%%%%%%%BIBTLATEX
\usepackage[
backend=biber,
style=numeric-comp,
natbib=true,
url=false, 
doi=true,
eprint=true,        
isbn=false,         
maxcitenames=2,     
maxbibnames=99,     
giveninits=true,
sorting=nty
]{biblatex}
\renewcommand*{\bibfont}{\small}
% Definición del encabezado personalizado con la línea
\usepackage{xurl}
% --- 7. DISEÑO EDITORIAL (Títulos, Índices, Fancyhdr) ---
\usepackage{epigraph}
\usepackage[margin=1cm, font=small, labelfont=bf, labelsep=endash]{caption}
\usepackage{booktabs}